\Needspace{8\baselineskip}\originalchapter{极向场设计与线圈反问题}\label{ch:pf}
\originalsection{固定边界与自由边界}
固定边界 GS 在指定闭合边界给 \ensuremath{\Psi}, 求内部平衡. 自由边界问题还要计算等离子体外真空场、PF 线圈和等离子体位置, 边界本身未知. 内部压力和 F 剖面、外部线圈、电源约束相互耦合. 不能把固定边界的数学解自动称作可由现实线圈维持的装置.

\originalsection{由等离子体电流求外场}
轴对称电流环的磁通用 Biot--Savart Green 函数 G 写成 $\Psi_P(R,Z)=\int G(R,Z;R',Z')J_\varphi(R',Z')\dd R'\dd Z'$. 为避免轴上或源点奇性, 数值积分要处理局部可积对数核. 线圈 c 的贡献是 $I_cG_c$, 因线性 Maxwell 真空方程, 线圈电流在给定位置时线性进入. 等离子体响应却可能非线性, 因其 J 与 \ensuremath{\Psi} 相依.

选边界采样点 j, 要求磁通等值或法向场满足目标, 得 $A_{jc}I_c\simeq b_j$. 方程通常过定或病态; 最小二乘 \begin{equation}\min_I\bigl(\|W(AI-b)\|^2+\lambda\|LI\|^2\bigr).\end{equation} 可约束电流、平衡平滑度或功耗. 用 QR/SVD 求解比直接平方成正规方程可靠, 且要报告奇异值、残差和最大线圈电流. 小残差但数千倍电流仍是不可实施设计.

\originalsection{真实 tokamak 与替代构型}
课程客座讲座介绍 tokamak 测量、控制及 stellarator、reversed-field pinch、spheromak 等替代构型. 共同问题是磁面、约束和 MHD 自由能, 不同点是环向/极向磁场由外线圈还是等离子体电流产生、对称性、变换和压力上限. 这些2007年的装置案例是历史教学背景; 本册不把预测建造时间写作2026事实, 也不重绘第三方装置照片. 理解平衡方法后, 读者可把各构型看作不同边界与通量函数问题, 而非记忆名称列表.
